\documentclass[class=scrbook, crop=false]{standalone}
\usepackage[subpreambles=true]{standalone}
\ifstandalone
    \input{../settings+/settings}
\fi

% ----------------------------------------------------------------------------
%                               Conclusion
% ----------------------------------------------------------------------------
\begin{document}

\chapter{Conclusion and Outlook} % Overview text
\label{Chapter::Conclusion}

This thesis regularized the indefinite QPs that arise in exact-Hessian SQP for NMPC at the QP level, inside the interior point solver,
where the barrier curvature is already taken into account and only the reduced Hessian needs to be positive definite.
We showed that the Riccati backward pass is a block $\mathrm{LD}\mathrm{L}^\top$ factorization of the reduced Hessian,
so that indefiniteness is detected and corrected stage-wise wherever the Cholesky factorization of $\tilde{M}_k$ fails,
while the linear complexity of the recursion in the horizon length is preserved.
We further showed that the minimum-perturbation rule propagates negative curvature through the cost-to-go matrix to the preceding stages.
This led to an improved strategy that uses a non-minimal correction and also regularizes the cost-to-go matrix,
and we proved that the regularization becomes inactive close to the solution, so that the local convergence rate of the interior point method is retained.

On a benchmark of $175$ indefinite QPs from eight OCPs,
whose analysis showed that the cost-to-go matrix is involved in most indefinite stages,
\textsf{PM\_MIRROR} left $4$ instances unsolved and \textsf{PM\_GMW} $7$ at a lower cost,
against $32$ to $59$ for full Hessian regularization and $15$ for inertia correction.
The full Hessian methods failed on almost every \textsl{Flying robot} instance, where the indefiniteness is confined to the cost-to-go matrix at a few stages,
and the instances left unsolved by all methods combine deep negative curvature with a large share of indefinite stages.

Several questions remain open:
\begin{itemize}
    \item \textbf{Regularization magnitude.} The eigenvalue floor $\epsilon$ is fixed,
    and the remaining failures occur where large corrections are applied at many stages.
    Adapting the magnitude, e.g.\ to the conditioning of the corrected blocks or to the barrier parameter, could improve robustness.
    \item \textbf{Other MPC formulations.} In time-optimal and economic MPC,
    the cost provides little positive curvature, so indefiniteness is likely to be broader than in this benchmark.
    Whether the ranking of the methods carries over remains to be tested.
    \item \textbf{Embedding in SQP.} The QPs were solved in isolation with a Python prototype.
    An implementation in HPIPM and \texttt{acados} is needed to compare against SQP-level regularization
    in terms of SQP iterations, run time and closed-loop performance.
\end{itemize}

\end{document}
